\subsection*{一步一步建立直观}

\paragraph{1. 求解就是追问输出来自哪里}

固定矩阵后，拖动左图输入并追踪右图红点。方程 $Ax=b$ 要求找到落到指定输出 $b$ 的输入。对于非零行列式的方阵，每个目标恰有一个原像；逆矩阵记录这条反向寻找规则。



\paragraph{2. 压扁时同时出现不可达与不唯一}

两列共线后，输出被限制在一条线上。线外目标没有解；线内目标通常有整条输入线作为原像。沿着被压成零的方向移动输入，输出不变，这就是核。不要把不可逆简单理解成每个方程都无解。



\paragraph{3. 把图像整理成解集的证明}

若 $x_0$ 已是一个解，且 $Av=0$，线性性给出 $A(x_0+v)=b$。反过来，任意另一个解减去 $x_0$ 后都被送到零。因此解集恰为 $x_0$ 加上核；这些平行原像族可以由这两步等式完整解释。保持压扁情形，沿核方向拖动输入，就能检查输出是否保持不变。
